\section{Síntesis y ampliación}

Esta sección condensa el episodio completo en un marco de estudio. Primero, los paradigmas de modelos no son artículos aislados: son el resultado de que el hardware, los datos, la arquitectura y las técnicas de entrenamiento se alineen. Segundo, la Infra y los datos amplifican la capacidad de los modelos. Sin fragmentación de la memoria de GPU, sin eficiencia de clúster y sin datos de alta calidad, la Scaling Law se queda en una regularidad sobre el papel. Tercero, el desarrollo de los modelos de lenguaje avanza de la representación al preentrenamiento y después al posentrenamiento y la alineación. Cuarto, el desarrollo multimodal es una convergencia gradual de la comprensión visual, la generación de imágenes, la alineación entre imagen y texto y la adopción del Transformer. Quinto, lo esencial al leer artículos no es terminarlos, sino desarrollar conciencia de sus límites.

\begin{importantbox}{Ubicar EP117 en la serie de IA de Zhang Xiaojun (张小珺)}
EP119 trata la arqueología de la arquitectura Attention. EP118 explica cómo entiende un CEO la VLA, el Agent OS y las plataformas. EP117 lleva estas preguntas más atrás en la historia: por qué ganó el Transformer, por qué la Infra y los datos se volvieron la línea principal y por qué los modelos de lenguaje y los multimodales terminaron por converger.
\end{importantbox}

\subsection{Tabla condensada de los 36 artículos}

\begin{longtable}{p{0.22\textwidth}p{0.30\textwidth}p{0.38\textwidth}}
\toprule
Línea principal & Artículo/hito & Importancia didáctica \\
\midrule
Paradigmas de modelos & Brook, AlexNet, ResNet & La GPU, el aprendizaje profundo y las redes residuales cambiaron juntos el límite de lo que se puede entrenar. \\
Paradigmas de modelos & seq2seq, Attention, Transformer & El modelado de secuencias pasa de comprimir el estado a modelar directamente las relaciones entre tokens. \\
Paradigmas de modelos & AlphaGo Zero, MoE, CoT, LoRA, ReAct & El aprendizaje por refuerzo, el enrutamiento entre expertos, las indicaciones para razonar, la adaptación de bajo costo y los Agents forman las capacidades de los sistemas posteriores. \\
Paradigmas de modelos & The Bitter Lesson & A largo plazo, el cómputo general y el aprendizaje a escala superan una y otra vez a las estructuras diseñadas a mano. \\
Infra/datos & ZeRO, MegaScale & La fragmentación y la ingeniería de clústeres determinan si un modelo grande puede entrenarse. \\
Infra/datos & Scaling Law, Chinchilla & Los parámetros, los datos y el presupuesto de cómputo deben guardar proporción. \\
Infra/datos & LAION-5B, RefinedWeb & Los datos abiertos y la ingeniería de datos web determinan el límite del corpus entrenable. \\
Modelos de lenguaje & Word2Vec, Google Translate & La representación de palabras y el despliegue en línea de modelos neuronales son la prehistoria de los modelos de lenguaje. \\
Modelos de lenguaje & GPT-1, BERT, GPT-2, GPT-3 & El paradigma de preentrenamiento pasa de las tareas de comprensión a la generación y a la emergencia de capacidades por escala. \\
Modelos de lenguaje & InstructGPT, Tulu 3 & El posentrenamiento, la alineación con instrucciones y la reproducción abierta hacen que los modelos sean utilizables. \\
Multimodal & DeepVideo, redes de dos flujos & La comprensión de video lleva el aprendizaje profundo de la imagen estática a la dimensión temporal. \\
Multimodal & GAN, Diffusion, DDPM & Los modelos generativos pasan del entrenamiento adversarial a la generación estable por difusión. \\
Multimodal & ViT, CLIP, Stable Diffusion, DiT & El Transformer visual, la alineación entre imagen y texto y el Transformer de difusión convergen. \\
\bottomrule
\end{longtable}

\subsection{Lecturas adicionales}

\begin{itemize}
\item Si le interesan los paradigmas de modelos, puede contrastar este episodio con la revisión de la arquitectura Attention de EP119, para entender por qué, después del Transformer, todavía se sigue «puliendo la Attention».
\item Si le interesa la IA del mundo físico, puede contrastarlo con EP118, EP120 y EP121, para reunir la VLA, los modelos del mundo y los datos de robots en un mismo mapa.
\item Si le interesan los Agents, puede contrastarlo con la historia técnica de los Agents de EP139 y con el paradigma de posentrenamiento/Agent de EP138, para entender cómo CoT/ReAct se convierten en sistemas reales.
\end{itemize}

\end{document}
